\section{The six-command sphere and variable accuracy}\label{sec:noidle62}
The full spectral measure in Lemma~\ref{lem:lpsgap61} already lives on
six nonidentity commands. The following result uses exactly those six
commands for the experiment itself, with no idle instruction.

\begin{theorem}[An explicit no-idle experiment]\label{thm:noidle62}
Let $A_1,A_2,A_3$ be the rational matrices in \eqref{eq:lpsinput61}.
Take exactly
\[
 \mathcal A_6=\{A_1,A_1^{\mathsf T},A_2,A_2^{\mathsf T},A_3,A_3^{\mathsf T}\},
 \qquad x_0=e_1,\qquad F(w)=D(A_we_1).
\]
Every decoder must lie in $\mathcal D_2$, and error is Frobenius error.
For $0<\varepsilon<1/\sqrt2$, put
$\zeta=1-\sqrt2\varepsilon$ and $\theta=(1-\zeta)/\zeta$.
Then every admissible width profile satisfies
\begin{equation}\label{eq:noidleoccupation62}
 \#\{1\le t\le N:|S_t|\le k\}
             \le1+6\log(12k)+648\pi^2\theta k.
\end{equation}
For $N\ge2$, the minimum peak width obeys the explicit joint bounds
\begin{equation}\label{eq:noidlejoint62}
 \min\left\{\frac{e^{(N-1)/12}}{12},
                   \frac{N-1}{1296\pi^2\theta}\right\}
 \le W_{N,\varepsilon}
 \le\min\{5^N,\ 52+104N/\varepsilon\}.
\end{equation}
At every horizon the exact simultaneous cut profile is $5^t$,
$0\le t\le N$, and the same profile is necessary and attained for
$0\le\varepsilon\le25^{-N}/16$. At
$\varepsilon\ge1/\sqrt2$, one label suffices.
For any fixed $\varepsilon_0<1/\sqrt2$ and sequence
$0<\varepsilon_N\le\varepsilon_0$ with
$\log(1/\varepsilon_N)=o(N)$,
\begin{equation}\label{eq:sixjointmatched62}
                  W_{N,\varepsilon_N}=\Theta(N/\varepsilon_N).
\end{equation}
The constants in this last comparison can be taken independent of the
sequence; its starting horizon can depend on the sequence.
\end{theorem}
\begin{proof}
The classical full spherical norm is $\sqrt5/3$ for the uniform law on
exactly $\mathcal A_6$ \cite[Theorem~1.1]{PLP61pre}. Use the drift proof
of Theorem~\ref{thm:driftoccupation62}, taking $h=k^{-1/2}$ and the
sharper sphere constants $\gamma=1/3$, $L_h=\log(12k)$, $B_3=24$
as in Theorem~\ref{thm:effectivelps61}. This gives
\eqref{eq:noidleoccupation62}, with arbitrary deterministic commands
in the gaps. It imports no uncomputed spectral constant.

For a full width bound $k$, at least one of
$6\log(12k)$ and $648\pi^2\theta k$ is at least $(N-1)/2$.
Inverting the corresponding inequality gives the lower bound in
\eqref{eq:noidlejoint62}. The sparse rational construction in
Theorem~\ref{thm:effectivelps61} applies unchanged to the six-letter
sublist; no row for the identity is needed. It gives the second upper
bound, with the same rational-decoder and dyadic-row guarantees for
rational tolerances.

For the exact profile, recall the independently proved free-group and
stabilizer calculation in Theorem~\ref{thm:effectivelps61}. The orbit is
the right-coset space of $\langle a\rangle$ in the free group on $a,b,c$,
where $a=A_1$. Every coset has a unique shortest representative $w$
with no terminal $a^{\pm1}$. If its length is $\ell\le t$, the group
word $w a^{t-\ell}$ has length $t$ and the same action on $e_1$.
In chronological execution this means applying the $a$'s first; they
fix the seed. There is no cancellation when $w$ is nonempty, and the
empty representative is padded by $a^t$. Thus even without an identity
letter, the length-$t$ reachable orbit is exactly the radius-$t$ coset
ball. Its size is $5^t$. The extreme-orbit theorem gives the exact
profile and deterministic orbit storage gives the upper bound $5^N$.

The denominator-five lattice and fixed-suffix isometry argument of
Theorem~\ref{thm:effectivelps61} does not require identity padding once
the orbit count is established. It gives the same robust interval.
The center $I/2$ gives the final constant-width assertion.

Finally $\theta\le C_{\varepsilon_0}\varepsilon_N$, while
$\log(N/\varepsilon_N)=o(N)$. The exponential alternative in the
minimum on the left of \eqref{eq:noidlejoint62} eventually exceeds a
constant times $N/\varepsilon_N$. The other alternative already has
that lower order. The upper bound proves \eqref{eq:sixjointmatched62}.
\end{proof}

\begin{remark}[Three separate accuracy regimes]\label{rem:threeaccuracy62}
The exact count and its robust interval, the explicit two-parameter
bounds \eqref{eq:noidlejoint62}, and the subexponential-accuracy
asymptotic \eqref{eq:sixjointmatched62} are different statements.
For example, $\varepsilon_N=N^{-a}$ gives width $\Theta(N^{a+1})$,
and $\varepsilon_N=e^{-N^b}$, $0<b<1$, gives
$\Theta(Ne^{N^b})$. Neither assertion extrapolates the exact plateau
to a fixed positive error or resolves its full crossover at
$\log(1/\varepsilon_N)\asymp N$. The constant $6500$ remains a safe
fixed-error simplification of \eqref{eq:noidleoccupation62}, not an
optimal leading constant. The seven-letter experiment and the older
rational-unitary pair remain separate retained examples.
\end{remark}
